\chapter{Programming Documentation}
\label{app:program}

The codebase can be divided into six main components:

\begin{itemize}
    \item \texttt{src/App} contains the main application logic and serves as the entry point of the system,
    
    \item \texttt{src/Evolution} implements the evolutionary algorithms, including the evolution of the packing rules vector and the evaluation of solution fitness,
    
    \item \texttt{src/Packing} handles the packing process, transforming the packing input and packing rules into a concrete solution,
    
    \item \texttt{src/SettingForm} provides the user interface for configuring the algorithms and their parameters,
    
    \item \texttt{tests/} contains unit tests for the C\# implementation,
    
    \item \texttt{scripts/} includes Python scripts for dataset generation, experimental evaluation, and result processing.
\end{itemize}

\section{\texttt{src/Packing}}

We first take a look at the packing module.

\texttt{PackingInput.cs} represents the input data for a packing problem
and contains container properties as well as the list of boxes to be packed.

\texttt{PackingSetting.cs} defines configuration settings controlling the behaviour
of the packing solver. It includes allowed placement heuristics, rotation permissions,
and an optional packing order heuristic.

\texttt{PackingSolver.cs} is the main entry point for solving a bin packing problem
using a set of packing rules. An instance of this class is provided with \texttt{PackingRules}
and \texttt{PackingSetting}. When its main function is called with a set of packing rules,
it first decodes the rules into a sequence of boxes and subsequently assigns them to containers.

\subsection{\texttt{src/Packing/Rules}}

This folder contains \texttt{PackingRules.cs} and \texttt{PackingRulesCell.cs},
which together form a list-based representation of the packing-rules vector.

The folder \texttt{Decoding} contains \texttt{Decoding/PackingRulesDecoder.cs},
which decodes packing rules into an ordered list of pending boxes.

The folder \texttt{Decoding/PartDecoders} contains additional interfaces and classes used for decoding
rotation and placement separately, while \texttt{Decoding/Sorting} contains the infrastructure
for sorting the pending boxes either by evolutionary order or by order heuristics, which are also defined in this folder
as described in Section~\ref{sec:order-heuristics}.

\subsection{\texttt{src/Packing/Architecture}}

This folder contains the data structures and helper classes related to boxes, containers, and the packing process in general. 
The main file is \texttt{BoxPacker.cs}, which packs pending boxes into containers according to the algorithm defined in Algorithm~\ref{alg:box-packing}.

The \texttt{Boxes} folder contains the definitions of box-related classes, including box properties, pending boxes, and placed boxes. 
In particular, these are \texttt{Boxes/BoxProperties.cs} (which is part of \texttt{PackingInput}), 
\texttt{Boxes/\\PendingBox.cs}, and \texttt{Boxes/PlacedBox.cs}, together with the corresponding box interfaces and extension methods.

The \texttt{Geometry} folder contains classes related to the geometric representation of the problem, including coordinates, sizes, regions, and rotations, as described in Chapter~\ref{chap:problem}. 
These are defined in \texttt{Geometry/Coordinates.cs}, \texttt{Geometry/\allowbreak Sizes.cs}, \texttt{Geometry/Region.cs}, and \texttt{Geometry/Rotation.cs}.

The \texttt{EMR} folder contains the implementation of empty maximal regions and the corresponding update algorithm.

The \texttt{Container} folder defines the opened container class \texttt{Container/\\OpenedContainer.cs}, 
which provides functionality for packing pending boxes into a container, as well as \texttt{Container/ContainerProperties.cs}, 
which is part of \texttt{PackingInput}. It also contains \texttt{Container/ContainerData.cs}, a data class used 
as one of the inputs to the placement heuristic, providing it with all relevant information about the container.

Placement is implemented in \texttt{Placements}. In particular, \texttt{Placements/\\Placement.cs} 
stores the container in which a box is placed and the region it occupies, while \texttt{Placements/PlacementHeuristic.cs} 
defines the interface for the placement heuristic function. The file \texttt{Placements/PlacementHeuristics.cs} 
contains the implementations of the placement heuristics described in Section~\ref{sec:placement-heuristics}.

\section{\texttt{src/Evolution}}

This section describes the implementation of the evolutionary algorithms and the related supporting classes.

We first look at the configuration classes in \texttt{Setting}. 
The main application configuration is defined in \texttt{Setting/ProgramSetting.cs}. 
It contains the path to the input file, the output file for the results, 
an additional file for storing statistics about the evolutionary run, 
as well as the packing settings and the evolutionary algorithm settings, which are described below.

The settings are defined in \texttt{Setting/EvolutionaryAlgorithmSettings}. 
This folder contains the base class \texttt{Setting/EvolutionaryAlgorithmSettings/\\EvolutionaryAlgorithmSetting.cs}, 
together with two derived classes corresponding to the implemented algorithms, 
namely \texttt{Setting/\\EvolutionaryAlgorithmSettings/ElitistGeneticAlgorithmSetting.cs} 
and \texttt{Setting/EvolutionaryAlgorithmSettings/\\ProbabilisticHillClimbingSetting.cs}. 
These classes define the hyperparameters of each algorithm. 
The base class contains parameters common to both algorithms, such as the population size and the number of generations, 
while the derived classes define algorithm-specific hyperparameters.

The class \texttt{EvolutionaryAlgorithms.cs} acts as a factory for creating evolutionary algorithms from the evolutionary algorithm settings. 
Depending on the provided setting instance, it returns the corresponding evolutionary algorithm configured with the specified hyperparameters. 
Its input also includes a fitness evaluator, which is used throughout the evolutionary process, 
a population factory, and a class for recording statistics about the run.

This statistics class is located in the folder \texttt{EvolutionStatistics} and is responsible for storing
a snapshot of each generation with relevant information such as the best fitness, average fitness, 
and the time required to progress to the next generation.
Each snapshot is represented by the class \texttt{EvolutionStatistics/\allowbreak StatisticalData.cs}, 
while the management of the full sequence of snapshots is defined by 
\texttt{EvolutionStatistics/\allowbreak EvolutionStatistics.cs}.

Regarding fitness evaluation, the folder \texttt{Fitness} contains the implementation of the fitness computation.
First, \texttt{Fitness/\allowbreak ContainerFitnessEvaluator.cs} evaluates the fitness of a 
given set of containers according to the formula defined in Section~\ref{sec:fitness-function}.
Second, \texttt{Fitness/\allowbreak PackingRulesFitnessEvaluator.cs} is responsible for evaluating packing rules vectors. 
It first uses \texttt{src/\allowbreak Packing/PackingSolver.cs} to obtain a packing solution, 
and then passes this solution to \texttt{Fitness/\\ContainerFitnessEvaluator.cs} to compute the final fitness value.
This class also supports evaluating an entire population of packing rules in parallel, which significantly reduces computation time.

The main file in \texttt{src/Evolution} is \texttt{EvolutionProgram.cs}, 
which takes the program settings, creates the fitness evaluator and the population factory, 
constructs the evolutionary algorithm with the given hyperparameters, 
executes the evolutionary process, and returns the best individual together with the statistics of the evolution.

\subsection{\texttt{src/Evolution/Architecture}}

This folder contains the core architecture used to build the evolutionary algorithms.
The main interface is \texttt{IEvolutionary.cs}. 
The class\\\texttt{EvolutionaryBase.cs} serves as a base class for evolutionary algorithms and implements the overall evolutionary process. 
For a specific derived algorithm, it is only necessary to implement the function that creates the next generation.
The algorithms are generic, so the individual can be represented by any type. 
The folder also contains the helper class \texttt{EvaluatedIndividual.cs}, which combines a generic individual with its fitness value.

The folder \texttt{ElitistGeneticAlgorithm} contains the implementation of the elitist genetic algorithm, as described in Algorithm~\ref{alg:ega-pseudocode}. 
It provides two variants: a generic implementation \texttt{ElitistGeneticAlgorithm/\allowbreak ElitistGenetic.cs}, 
and a wrapper specialized for packing rules, \texttt{ElitistGeneticAlgorithm/\\PackingRulesElitistGenetic.cs}.

Similarly, the folder \texttt{ProbabilisticHillClimbing} contains the implementation of probabilistic hill climbing, as described in Algorithm~\ref{alg:phc-pseudocode}. 
It also includes two variants: a generic implementation \texttt{ProbabilisticHillClimbing/\allowbreak ProbabilisticHillClimbing.cs}, 
and a packing rules specific wrapper\\\texttt{ProbabilisticHillClimbing/PackingRulesProbabilisticHillClimbing.cs}.

The folder \texttt{Population} contains \texttt{Population/\\PackingRulesPopulationFactory.cs}, 
a class that receives the length of the packing rules vector upon initialization and provides a function for generating a specified number of random individuals.

The remaining parts of the architecture consist of various evolutionary components, namely the folders 
\texttt{Mutation}, \texttt{Selection},
\texttt{Crossover}, \texttt{Elitism}, and 
\texttt{Memetic}. These folders generally follow a consistent structure: each contains a generic interface for the corresponding component, 
together with concrete implementations for packing rules.

This design is useful because the evolutionary algorithms depend only on these components through their interfaces, making it easy to replace individual operators. 
For example, in \texttt{Selection}, both random selection and tournament selection are implemented, and both conform to the same interface.

\section{\texttt{src/SettingForm}}

This folder contains the Windows Forms implementation for manually selecting the algorithm and its parameters. 
The main entry point is \texttt{FormProgram.cs}, which first opens \texttt{SettingForm.cs}. 
In this form, the user selects the algorithm name, input and output files, and the packing settings.

After this step, the program selects the form from a dictionary based on the chosen algorithm, either 
\texttt{ProbabilisticHillClimbingSettingForm.cs} or 
\texttt{ElitistGeneticSettingForm.cs}. 
In these forms, the hyperparameters of the evolutionary algorithm are specified.

Once the configuration is completed, the function returns a new program setting based on the values provided in the forms.

\section{\texttt{src/App}}

This folder serves as the entry point of the system.
The main function in \texttt{Program.cs} either takes a command-line argument 
with the path to the JSON file containing the program settings, or launches the Windows Forms 
application described above if no argument is provided. If a path is given, 
the class \texttt{Input/ProgramSettingLoader.cs} is used to load the program settings 
from the JSON file. Once this is done, the class \texttt{Input/PackingInputLoader.cs} 
is used to load the packing input from the path specified in the program settings.

The class \texttt{src/Evolution/EvolutionProgram.cs} is then used together with 
the loaded settings and packing input to run the program and obtain the results, 
namely the best individual and the statistics of the evolutionary run.

Both of these outputs are then saved using helper classes from the folder \texttt{Output}. 
In particular, \texttt{Output/EvolutionStatisticsSaver.cs} saves the statistics as a CSV file, 
while \texttt{Output/PackingOutputSaver.cs} saves the best solution as a JSON file.

\section{\texttt{scripts}}

This folder contains the Python scripts used for dataset generation, evaluation, and result visualization.

The script \texttt{generate\_full\_mpv\_benchmark.py} was used to generate the full dataset described in Section~\ref{sec:datasets}. 
It relies on several helper scripts:\\ 
\texttt{generate\_mpv\_benchmark.py}, which generates a dataset with a given mean weight, standard deviation, number of instances, and related parameters; 
\\\texttt{mpv\_generator.py}, which creates a single packing instance for a given class and number of boxes using the same parameters; 
\texttt{saver.py}, which saves the packing input as a JSON file; and 
\texttt{models.py}, which provides Python versions of the necessary data structures, such as container and box properties.

Once the dataset has been generated, \texttt{lower\_bound.py} can be used to compute the average lower-bound fitness for each class and number of boxes, as shown in Table~\ref{tab:lower-bound}.

The script \texttt{run\_experiments.py} is used to run the evolutionary algorithm with given hyperparameters on the full dataset. 
It takes the desired evolutionary and packing settings, together with the C\# solver, and creates temporary JSON files containing the corresponding program settings. 
It then iterates over all instances in the dataset, runs the solver, and stores both the results and the evolutionary statistics for each instance. 
In addition, it computes average fitness values per generation across all instances of the same class and number of boxes.

These average statistics can then be processed by \texttt{plot.py}, which generates graphs for the average results for each class and number of boxes. 
The script can also combine the results of multiple evolutionary algorithms and plot them in the same figure.